Relation (48.2) also admits another interpretation. The integral
$\oint p\,dx$ is the area enclosed by the particle's closed

\SourcePage{215}{218}
classical phase trajectory (that is, by the curve in the $p,x$ plane---the
particle's phase space). If this area is partitioned into cells, each of area
$2\pi\hbar$, their total number is $n$. But $n$ is the number of quantum
states whose energies do not exceed the specified value (the value
corresponding to the phase trajectory under consideration). Thus, in the
quasiclassical case, one may associate with every quantum state a phase-space
cell of area $2\pi\hbar$. Equivalently, the number of states assigned to a
phase-space area element $\Delta p\Delta x$ is
\begin{equation}
  \frac{\Delta p\Delta x}{2\pi\hbar}.
  \tag{48.4}
\end{equation}
If the wave vector $k=p/\hbar$ is used in place of the momentum, this number
takes the form $\Delta k\Delta x/2\pi$. As expected, it agrees with the known
expression for the number of normal modes of a wave field (see II, §~52).
